\chapter{引言}
\label{chap:introduction}

\section{研究背景与动机}

图像语义分割是计算机视觉领域的核心任务之一，旨在对图像中的每个像素进行分类。与图像分类仅输出单一标签不同，语义分割需要生成与输入图像等尺寸的稠密预测图，这对模型的空间感知能力提出了更高要求。

本实验选择 Oxford-IIIT Pet Dataset 作为研究对象，该数据集包含 37 种宠物的图像及其三分类标注：前景、背景和轮廓。该任务面临三方面挑战：其一，轮廓像素占比极低，不足 5\%，严重的类别不平衡使模型倾向于忽略边界区域；其二，宠物属于非刚性物体，姿态多样，要求模型具备较强的泛化能力；其三，复杂的光照条件、物体遮挡以及前景与背景颜色相近等因素进一步增加了分割难度。

\section{研究目标}

本实验的核心目标包含三个层次。首先，构建一个可复现的 baseline 模型，采用手写 U-Net 从零训练，建立性能基准。其次，通过数据增强、迁移学习、损失函数设计、训练策略调优等多层次优化手段，系统性地提升分割性能，将 mIoU 从 0.71 提升至 0.81，实现 10.5\% 的提升。最后，通过消融实验量化各优化项的独立贡献与协同效应，为类似分割任务提供可参考的优化路径。

\section{报告结构}

本报告共分为六章：第二章介绍数据集与评估指标；第三章构建 baseline 模型；第四章进行多层次优化实验；第五章展示实验结果与分析；第六章总结全文并展望未来工作。
